\subsection{Synthesis: Deployment Context}
\label{sec:analysis-deployment-synthesis}

Stage~K and Stage~C show that the architectural pattern remains visible after the
workload is moved into Kubernetes-style service chains. In local
\texttt{kind}, latency increases from the monolithic Kubernetes baseline
through the four-service chain, while the final five-service increment is
small because the added late-stage boundary transfers a much smaller
activation tensor. The same ordering and the same bounded-non-linearity
pattern also survive the kind multi-worker CNI sensitivity probe of
the Stage~K cross-node alternative (\cref{sec:analysis-rq14b}), so neither finding is an artefact
of single-node pod colocation on the local cluster; the kind multi-node
per-condition deltas are, however, an order of magnitude larger than the
AKS multi-node deltas of \cref{sec:analysis-rq15b} and should be read as
a kind-CNI sensitivity envelope, not as a substitute for a managed-cloud
cross-host measurement. In AKS, the same condition ordering is preserved:
\texttt{monolithic\_k8s\_1svc}, \texttt{chain\_2svc}, \texttt{chain\_3svc},
\texttt{chain\_4svc}, and \texttt{chain\_5svc}. The absolute millisecond
values differ between local Kubernetes and AKS, but the within-environment
overhead trend transfers directionally, as summarised in
\cref{tab:rq15_transfer}.

This distinction is important. The thesis does not claim that local Kubernetes
predicts AKS latency numerically. Cloud and Kubernetes performance studies
show that container runtime, virtualization, CNI, managed-cluster
implementation, and cloud variability can change absolute
measurements~\cite{managed_k8s_perf,virtualization_costs,cni_k8s_ic2e21,
noise_in_clouds}. The stronger supported claim is architectural: when the
same ResNet-18 chain family is run in AKS, the relative ordering remains
interpretable even though the platform changes.
